\documentclass[12pt,a4paper]{article}
% Preambule commun aux comptes rendus CR_*.tex.

\usepackage{array}
\newcolumntype{C}[1]{>{\centering\arraybackslash}m{#1}}
\usepackage{multirow}
\usepackage{geometry}
\geometry{a4paper, margin=1in}
\usepackage{graphicx}
\usepackage{float}
\usepackage{tikz}
\usetikzlibrary{arrows.meta,positioning}
\usepackage{pgfgantt}
\usepackage{pdflscape}
\usepackage{enumitem}
\usepackage{booktabs}
\usepackage{longtable}
\usepackage{ragged2e}
\usepackage{tabularx}
\usepackage{xparse}
\usepackage{ltablex}
\usepackage{xltabular}
\usepackage{subcaption}

\newcolumntype{Y}{>{\RaggedRight\arraybackslash}X}

\newcommand{\StyledTableSetup}{%
  \footnotesize
  \renewcommand{\arraystretch}{1.2}%
  \setlength{\tabcolsep}{4pt}%
}


% ============================================================
% IHEDN COVER TEMPLATE (XeLaTeX)
% ============================================================
% Compilation (from project root):
%   latexmk -pdf main.tex
%   LATEXMK_SHELL_ESCAPE=0 latexmk -pdf main.tex
%
% Optional environment controls from latexmkrc:
%   LATEXMK_ENGINE=xelatex        (default/enforced)
%   LATEXMK_SHELL_ESCAPE=1|0      (default 1, needed for SVG conversion)
%
% Required package for this template:
%   \usepackage{ihedn-cover}
%
% Note: ihedn-cover already loads needed internal packages
% (fontspec, graphicx, tikz, ragged2e, optional svg).
% ============================================================

% ----------------------------------------------------------------
% Package options (documented)
% ----------------------------------------------------------------
% Geometry scope:
%   apply-geometry            -> force A4 + zero margins while cover is built
%   no-apply-geometry         -> do not alter host geometry (default)
%
% Typesetting freeze:
%   global-typesetting        -> apply freeze globally (intrusive)
%   no-global-typesetting     -> no global freeze (default)
%
% Small-caps handling:
%   local-smallcaps-fix       -> \textsc -> \MakeUppercase only inside cover (default)
%   no-local-smallcaps-fix    -> disable local fix
%   global-smallcaps-fix      -> global override (intrusive)
%   no-global-smallcaps-fix   -> disable global fix (default)
%
% Main font handling:
%   global-mainfont           -> set Times New Roman globally (default)
%   local-mainfont            -> set Times only during cover rendering
%   no-mainfont-override      -> keep host document main font
%
% SVG support:
%   svg-support               -> enable SVG logos (default, needs shell-escape)
%   no-svg-support            -> disable SVG path handling (pdf/png only)
%
% Debug overlay:
%   debug                     -> show mm grid + bounding boxes on cover
%   no-debug                  -> hide debug overlay (default)
% ----------------------------------------------------------------
\usepackage[
  no-apply-geometry,
  no-global-typesetting,
  global-smallcaps-fix,
  local-smallcaps-fix,
  global-mainfont,
  svg-support,
  no-debug
]{ihedn-cover}

% ----------------------------------------------------------------
% Bibliography stack (BibLaTeX/Biber, style from subtree vendor/)
% ----------------------------------------------------------------
\usepackage{polyglossia}
\setdefaultlanguage{french}
\makeatletter
\g@addto@macro\captionsfrench{%
  \renewcommand{\tablename}{Tableau}%
  \renewcommand{\figurename}{Figure}%
}
\makeatother
\usepackage{csquotes}
\usepackage{xurl}
\usepackage{hyperref}
\makeatletter
\let\ihedn@orig@input@path\input@path
\def\input@path{{vendor/biblatex-sciences-po-fr/style/}}
\makeatother
\usepackage[
  backend=biber,
  style=sciencespo,
  hyperref=true,
  autolang=other
]{biblatex}
\AtBeginBibliography{\sloppy\setlength{\emergencystretch}{3em}}
\makeatletter
\let\input@path\ihedn@orig@input@path
\makeatother
\addbibresource{bibliographies/bib/rapport.bib}

% ----------------------------------------------------------------
% tcolorbox
% ----------------------------------------------------------------
\usepackage[most]{tcolorbox}
\tcbuselibrary{breakable}

% Example: use Arial for the full report + cover.
% Uncomment and switch package option to no-mainfont-override.
% \setmainfont{Arial}[
%   UprightFont    = *,
%   BoldFont       = * Bold,
%   ItalicFont     = * Italic,
%   BoldItalicFont = * Bold Italic
% ]

%==================================================
% Liste des propositions
%==================================================

\makeatletter

\newcommand{\listpropositionname}{Liste des propositions}

\newcommand{\listofpropositions}{
  \section*{\listpropositionname}
  \addcontentsline{toc}{section}{\listpropositionname}
  \@starttoc{lop}
}

\newcommand*\l@proposition{\@dottedtocline{1}{1.5em}{2.3em}}

\makeatother

%==================================================
% Compteur des propositions
%==================================================

\newcounter{proposition}

%==================================================
% Style graphique des propositions
%==================================================

\newtcolorbox{propositionbox}[2][]{
  lower separated=false,
  colback=white!80!gray,
  colframe=white!20!black,
  fonttitle=\bfseries,
  colbacktitle=white!30!gray,
  coltitle=black,
  enhanced,
  sharp corners,
  attach boxed title to top left={
    xshift=0.5cm,
    yshift=-2mm
  },
  title=#2,
  #1
}

%==================================================
% Environnement proposition
%==================================================

\newenvironment{proposition}[1]
{
  \refstepcounter{proposition}

  \addcontentsline{lop}{proposition}{%
    \protect\numberline{\theproposition}#1%
  }

  \begin{propositionbox}
  {Proposition~\theproposition~: #1}
}
{
  \end{propositionbox}
}

% Renommage Liste des figures en Table des figures
\addto\captionsfrench{%
  \renewcommand{\listfigurename}{Table des figures}
}

\newcommand{\heure}[2]{{\NoAutoSpacing #1:#2}}

\DeclareFieldFormat[misc]{title}{#1}

% Activer ce package avec l'option none retire la césure.
%\usepackage[none]{hyphenat}

\usepackage{adjustbox}
\usepackage{pdfpages}

\makeatletter
\newcommand{\IhednSetPdfPageCount}[2]{%
  \begingroup
  \edef\ihedn@pdfpath{#2}%
  \ifXeTeX
    \global#1=\XeTeXpdfpagecount "\ihedn@pdfpath" %
  \else
    \ifLuaTeX
      \saveimageresource{\ihedn@pdfpath}%
      \global#1=\lastsavedimageresourcepages\relax
    \else
      \pdfximage{\ihedn@pdfpath}%
      \global#1=\pdflastximagepages\relax
    \fi
  \fi
  \endgroup
}
\makeatother

\newcount\IhednIncludedPdfPageCount
\newcommand{\IhednRenderPdfPagesInBoxes}[1]{%
  \IhednSetPdfPageCount{\IhednIncludedPdfPageCount}{#1}%
  \foreach \IhednIncludedPdfPage in {1,...,\the\IhednIncludedPdfPageCount}{%
    \begin{tcolorbox}[
      colback=gray!5,
      colframe=black,
      boxrule=0.5pt,
      arc=3mm
    ]
    \centering
    \includegraphics[page=\IhednIncludedPdfPage,width=0.95\textwidth]{#1}
    \end{tcolorbox}
    \ifnum\IhednIncludedPdfPage<\IhednIncludedPdfPageCount
      \clearpage
    \fi
  }%
}

\usepackage{tablefootnote}


\begin{document}

% ----------------------------------------------------------------
% Cover keys (all available keys are shown below)
% ----------------------------------------------------------------
% Header block (top-right), logo, title, report lines, subtitle, committee.
\PageDeGardeIHEDN{
    header_1={Union-IHEDN},
    header_2={Association des auditeurs IHEDN},
    header_3={région Paris / Île de France},
    date={le 07 avril 2026},
    logo_path={Figures/logos_IHEDN/Logo_AR_16_PARIS_IDF.jpg},
    title={Crises majeures : quel rôle pour les citoyens engagés et les associations IHEDN ?},
    title_max_lines=2,
    title_fontsize={26.000pt},
    title_baselineskip={28.667pt},
    reportline_1={Rapport du Comité d'étude de l'Association régionale des auditeurs},
    reportline_2={de l'IHEDN région Paris / Île de France (AR 16)},
    subtitle={Compte rendu de la réunion du 19 Mars 2026},
    committee={Comité d'étude, d'octobre 2025 à juin 2026, composé de : Mme Valérie Arlé-Faivre, \textit{présidente} ; M. Aurélien Duvignac-Rosa, \textit{rapporteur} ; M. Pascal Bayot Duponchel ; M. Jean-Jacques Cleynen ; Mme Valérie Cowan ; M. Alain Fontan ; M. Pierre-Marie Giard ; Mme Muriel Joyeux ; M. Bruno Leray ; M. Yves-Henri Renhas ; M. Marc Roure ; Mme Isabelle de Segonzac ; M. Gérard Turck ; Mme Marie Danielle Vasquez-Duchêne.}
}

\section{Éléments de rappel}

La réunion du 19 mars 2026 s'inscrit dans la continuité des réflexions engagées
par le comité d'études autour du rôle potentiel des membres de l'AR~16 dans
le cadre de crises majeures.

L'un des premiers éléments portés à notre attention fut l'adhésion des anciens
auditeurs et auditeurs associés et leur implication dans les associations.
Cet élément s'inscrit dans la continuité de la réflexion abordée le 12 février
au sujet de l'annuaire lacunaire de l'IHEDN.

Par ailleurs, il a également été soulevé par les membres du comité que les adhérents
des associations IHEDN sont assez faibles par rapport aux membres d'autres associations
et ceux qui continuent activement le sont encore moins.
Il est primordial d'identifier les compétences collectives de ses membres (en s'appuyant
sur le questionnaire).

Dans ce sens, Gérard Turck (membre du CODIR), suite à un échange avec le bureau directeur,
nous a indiqué que la mise en place d'un questionnaire était une idée d'intérêt
pour le CODIR.

\subsection{Une volonté de mieux connaître et mobiliser les auditeurs}

Dans le prolongement des discussions du bureau, Monsieur Gérard Turck a évoqué la mise en place d'un questionnaire, ayant reçu un aval officieux, afin de mieux appréhender les profils, les attentes et les capacités d'engagement des membres. Cette démarche vise à cartographier les compétences mobilisables et à structurer une éventuelle contribution des auditeurs en situation de crise.

Alain Fontan a, pour sa part, enrichi la réflexion en rappelant plusieurs initiatives et travaux antérieurs menés au sein de l'Union IHEDN :
\begin{itemize}
    \item les réflexions historiques autour de la \textit{Revue Défense} et de l'annuaire des auditeurs ;
    \item la mission confiée à un étudiant, Ulysse Cordier, chargé de diffuser un questionnaire auprès des présidents d'associations régionales ;
    \item la perspective d'exploiter ces travaux, notamment via des échanges avec les membres du comité.
\end{itemize}

Il a également été souligné l'absence de travaux académiques structurants consacrés à l'IHEDN, ce qui conforte l'intérêt de capitaliser sur les initiatives existantes (albums d'auditeurs, colloques, productions historiques). À cet égard, le colloque organisé au Sénat en 2018 sur la résilience nationale et la place des citoyens a été cité comme référence, de même que plusieurs jalons historiques (1998, 2009, 2024).

Enfin, la proposition de Monsieur Fabrice Hamelin relative au développement d'un centre d'étude social de la défense a été mentionnée comme une piste de structuration complémentaire\footcite{hamelinIHEDN}.

\subsection{La question des réserves et de l'engagement opérationnel}

Les interventions de Bruno Leray ont apporté un éclairage opérationnel centré sur les dispositifs de réserve existants. S'appuyant sur un rapport de l'Inspection générale des affaires sociales (2016), il a été rappelé que la majorité des réservistes relève des sapeurs-pompiers, que les associations représentent un volume significatif d'environ 60\,000 personnes, et que d'autres formes de réserves existent (communales, sanitaires).

Un document du Sénat met en évidence l'efficacité de ces réserves, notamment en phase post-crise, tout en soulignant la nécessité de renforcer les dispositifs de préparation en amont.

Ces éléments conduisent à relativiser la capacité d'engagement immédiat des auditeurs IHEDN, tout en renforçant l'intérêt du questionnaire pour :
\begin{itemize}
    \item identifier les membres déjà engagés dans des dispositifs de réserve ;
    \item préciser leur niveau de disponibilité et leur champ d'action.
\end{itemize}

Dans ce contexte, plusieurs pistes ont été évoquées :
\begin{itemize}
    \item la création d'une « réserve stratégique » orientée vers les retours d'expérience (RETEX) ;
    \item la valorisation de l'appartenance à l'IHEDN comme label facilitant la reconnaissance institutionnelle.
\end{itemize}

\subsection{Vers une structuration de la contribution des auditeurs}

Au-delà des constats, les échanges traduisent une convergence vers un objectif commun : mieux structurer la contribution des auditeurs IHEDN, en conciliant réalisme (disponibilité, compétences, cadre d'action) et ambition (mobilisation et utilité opérationnelle).

Cette réflexion s'inscrit dans la volonté de préciser la plus-value de l'IHEDN dans le cycle de gestion des crises, en particulier en amont et en aval, en complémentarité des autorités compétentes.

\section{Éléments de réflexion sur la singularité et la contribution des auditeurs IHEDN}

Les échanges ont permis d'interroger la singularité des auditeurs de l'IHEDN dans la gestion des crises, ainsi que leur capacité à apporter une réelle plus-value au sein de dispositifs déjà structurés.

\subsection{Une valeur ajoutée fondée sur l'expertise et la transversalité}

Les membres ont souligné que la spécificité des auditeurs IHEDN réside dans la diversité de leurs parcours, leur expérience professionnelle et leur capacité à appréhender des problématiques complexes de manière transversale. Cette singularité pourrait se traduire par une contribution renforcée dans les domaines de la prospective, de l'analyse stratégique et de la production de scénarios.

Dans cette perspective, il a été évoqué que les auditeurs pourraient s'organiser afin de développer des travaux de prospective, fondés sur le partage d'expériences et d'expertises sectorielles, contribuant ainsi à éclairer la décision publique. Cette réflexion s'inscrit dans une problématique plus large de lien entre savoir et pouvoir, illustrée notamment par le recours à des conseils scientifiques au plus haut niveau de l'État.

\subsection{Un questionnaire structurant et sélectif}

Le questionnaire envisagé constitue un outil central de cette démarche. Il devra permettre non seulement d'identifier les compétences et les niveaux d'engagement des membres, mais également de filtrer progressivement les répondants en fonction de leur implication réelle.

Une attention particulière devra être portée à la « charge du répondant »\footnote{cf. entretien du 20/02/2026 avec le docteur Victor Violier, sociologue à l'IRSERM.}, afin d'éviter un abandon en cours de questionnaire, y compris de la part des profils les plus motivés. Il conviendra ainsi de trouver un équilibre entre exhaustivité et accessibilité, tout en intégrant des mécanismes incitatifs favorisant la complétude des réponses.

Par ailleurs, il a été proposé d'intégrer une question ouverte visant à identifier les membres appartenant à d'autres associations d'intérêt général, afin de mieux appréhender les réseaux d'engagement existants.

\subsection{Articulation avec les acteurs institutionnels et territoriaux}

Les échanges ont également porté sur les modalités d'interaction avec les acteurs institutionnels, notamment les délégués militaires départementaux (DMD)\footcite{DMD_def}. Il a été envisagé que les associations IHEDN puissent se mettre à leur disposition, notamment en proposant des actions de sensibilisation ou des conférences, contribuant ainsi à alléger certaines charges.

Dans cette optique, il a été rappelé que l'accès aux référents défense passe par les DMD, soulignant la nécessité d'inscrire toute démarche dans un cadre institutionnel clairement identifié.

\subsection{Axes de travail : RETEX, prospective et diffusion}

Deux axes principaux de contribution ont été identifiés :
\begin{itemize}
    \item le développement de travaux de retour d'expérience (RETEX) ;
    \item la conduite d'analyses prospectives.
\end{itemize}

La question de la diffusion de ces travaux a toutefois été soulevée. Il a été constaté que de nombreux rapports produits ne font pas toujours l'objet d'une valorisation suffisante, demeurant parfois sans suite malgré leur qualité. Dès lors, la nécessité d'améliorer les circuits de diffusion, notamment vers l'Institut et l'Union, a été soulignée.

Une réflexion a également été engagée sur les synergies possibles entre associations régionales, notamment entre territoires aux caractéristiques complémentaires (urbains et ruraux), afin d'enrichir les analyses produites. Dans le cas de l'AR~16, son rattachement à Paris conduit à privilégier une approche centrée sur des thématiques nationales.

\subsection{Place du citoyen et perspectives de valorisation}

Enfin, les échanges ont rappelé le rôle central des citoyens dans la gestion des crises, ceux-ci constituant les premiers acteurs présents sur le terrain. Cette réalité conforte l'importance de renforcer la résilience collective et de mieux intégrer cette dimension dans les travaux du comité.

Dans cette perspective, il a été proposé de produire un rapport d'étape substantiel, présentant les avancées des travaux et les principaux enseignements. Cette restitution, notamment en lien avec l'IRSEM, pourrait constituer un levier de mobilisation des membres, en les incitant à participer au questionnaire.

Enfin, l'hypothèse d'un prolongement des travaux du comité sur une année supplémentaire a été évoquée, afin de consolider les analyses et d'approfondir les propositions formulées.

%\newpage

%\printbibliography[heading=bibintoc,title={Références bibliographiques}]

%\nocite{*}

% ----------------------------------------------------------------
% Notes d'usage:
% - Le style bibliographique est dans:
%   vendor/biblatex-sciences-po-fr/style/
% - Ajouter des entrées dans les fichiers bibliographiques chargés ci-dessus:
%   vendor/biblatex-sciences-po-fr/bibliographies/{articles,livres,online,rapports,theses}.bib
% - Compilation requise:
%   latexmk -pdf main.tex
%   (latexmkrc est configuré pour XeLaTeX + Biber).
% ----------------------------------------------------------------

\end{document}
